% =============================================================
%  Fault zones and selectivity — zone identifiers, the two identifier
%  algorithms and the changeover to STATCOM operation. Figure: fault
%  zone allocation as inline TikZ. Principles; the recorded changeover
%  is in pugalur-thrissur.tex.
% =============================================================
\subsection{Fault Zones and Selectivity}
\label{sec:dcf-zones}

\subsubsection{Zone identification}
\label{sec:dcf-zones-identification}

A restart is permitted only for a fault on the overhead line; a fault in a
cable, a converter station or the AC network is permanent by nature or is
cleared by other protection. On a hybrid DC circuit the reliable
identification of the zone in which the fault occurred is therefore the
prerequisite for an automatic recovery~[24]. The protection of each link is
divided into the fault zones AC network, converter station, DC overhead
line and DC cable at each end, Figure~\ref{fig:dcf-zones}, with the HVDC
link zone spanning both converters and the DC circuit~[24], [25].

Voltage and current measurements at the zone boundaries --- current
transformers at the transition station separate the overhead line from the
cable --- produce fault zone identifiers that state on which side of the
measuring point the fault lies. A logic combines the identifiers from all
locations while the fault clearing is in progress, so that the zone is known
before the recovery sequence is released~[24]. Since the blocking of
Section~\ref{sec:dcf-faults-converter} does not depend on the zone, the
detection and the selectivity are separate functions with separate time
budgets: microseconds for the first, the fault clearing time for the
second.

Two identifier algorithms cover the range of fault impedances an overhead
line produces~[24]. Faults of low impedance are detected by a wave-front
logic that evaluates the polarity of the instantaneous changes of the direct
current. Faults of high impedance are detected by a residual-current logic
acting on the discharge currents of the DC cable; it extends the coverage to
fault impedances of several k$\Omega$.

\begin{figure}[!ht]
  \centering
  \begin{tikzpicture}[font=\footnotesize, >=Latex, line width=0.7pt]
    % AC networks
    \draw[pattern=north east lines, pattern color=atBorderStrong] (0,-0.5) rectangle (0.9,0.5);
    \draw[pattern=north east lines, pattern color=atBorderStrong] (12.1,-0.5) rectangle (13.0,0.5);
    % station A
    \draw (0.9,0) -- (1.6,0);
    \draw (1.6,0) -- (2.1,0.25);
    \draw (2.15,0) -- (2.6,0);
    \draw (2.85,0) circle (0.25); \draw (3.2,0) circle (0.25);
    \draw (3.45,0) -- (3.8,0);
    \draw (3.8,-0.6) rectangle (4.9,0.6); \draw (3.8,-0.6) -- (4.9,0.6);
    \node at (4.05,0.3) {$\approx$}; \node at (4.65,-0.3) {$=$};
    % DC circuit
    \draw (4.9,0.35) -- (8.1,0.35); \draw (4.9,-0.35) -- (8.1,-0.35);
    \draw[line width=2.2pt] (6.9,0.35) -- (8.1,0.35);
    \draw[line width=2.2pt] (6.9,-0.35) -- (8.1,-0.35);
    % station B
    \draw (8.1,-0.6) rectangle (9.2,0.6); \draw (8.1,-0.6) -- (9.2,0.6);
    \node at (8.35,0.3) {$\approx$}; \node at (8.95,-0.3) {$=$};
    \draw (9.2,0) -- (9.55,0);
    \draw (9.8,0) circle (0.25); \draw (10.15,0) circle (0.25);
    \draw (10.4,0) -- (10.9,0);
    \draw (10.9,0) -- (11.4,0.25);
    \draw (11.45,0) -- (12.1,0);
    % zone boundaries
    \foreach \x in {1.3, 5.6, 6.9, 8.6, 11.7} {
      \draw[atDanger, dashed, line width=0.6pt] (\x,-1.0) -- (\x,1.0);
    }
    % zone labels
    \draw[<->, atFgMuted] (0.0,-1.25) -- (1.3,-1.25) node[midway, below, text=atFg] {AC network A};
    \draw[<->, atFgMuted] (1.3,-1.25) -- (5.6,-1.25) node[midway, below, text=atFg] {converter station A};
    \draw[<->, atFgMuted] (5.6,-1.25) -- (6.9,-1.25) node[midway, below, text=atFg] {DC OHL};
    \draw[<->, atFgMuted] (6.9,-1.25) -- (8.6,-1.25) node[midway, below, text=atFg] {DC cable};
    \draw[<->, atFgMuted] (8.6,-1.25) -- (11.7,-1.25) node[midway, below, text=atFg] {converter station B};
    \draw[<->, atFgMuted] (11.7,-1.25) -- (13.0,-1.25) node[midway, below, text=atFg] {AC network B};
    \draw[<->, atFgMuted] (1.3,1.25) -- (11.7,1.25) node[midway, above, text=atFg] {HVDC link zone};
  \end{tikzpicture}
  \caption{Fault zones of one link with a hybrid DC circuit, after~[24],
           [25]; the thick section is the DC cable between the transition
           station and station~B}
  \label{fig:dcf-zones}
\end{figure}

\subsubsection{Changeover to STATCOM operation}
\label{sec:dcf-zones-statcom}

When a sequence ends unsuccessfully, Section~\ref{sec:dcf-sequence-classification},
both stations change over to STATCOM operation after the fault clearing:
the DC disconnect switches separate each converter from the DC circuit and
each station restarts independently~[24]. The changeover time is the sum of
the operating time of the slow DC disconnect switches and, when the link was
running at rated voltage, of the tap-changer movement to the STATCOM
position; with slow disconnects it ranges from about 15\,s to about a
minute~[24]. The converters are therefore unavailable for the transient
voltage support in the first seconds after the loss of transmission and
available for the steady-state support afterwards.

The same changeover is started directly, without any restart attempt, when
the fault zone is the DC cable~[24], [25]. A fault in one converter station
prohibits the changeover for that station, and only the other one
reconfigures to STATCOM operation~[24].

\FloatBarrier
